\ifja
\chapter{Closing（決算締め）の離散時間モデル}
\label{chap:closing_model}

簿記教育で神秘化される「損益振替」の正体は、物理系における\textbf{「フロー用レジスタのゼロクリア」と「ストックへの階差合流」}という離散時間アルゴリズムにすぎない。

\section{漸化式による純資産（Equity）の発展}
会計期間 $t \in \{0, 1, 2, \dots\}$ を離散時間ステップと定義する。損益計算書（P/L）の要素である収益 $\text{Rev}_t$ と費用 $\text{Exp}_t$ は、区間 $(t-1, t]$ における流量（Flow）である。

第 $t$ 期における純利益（Net Income）は階差として定義される。
\begin{equation}
    \Delta \text{Equity}_t \equiv \text{Rev}_t - \text{Exp}_t
\end{equation}
\begin{equation}
    \mathbf{\text{Equity}_t = \text{Equity}_{t-1} + \Delta \text{Equity}_t}
\end{equation}
\begin{equation}
    \text{Equity}_n = \text{Equity}_0 + \sum_{k=1}^n (\text{Rev}_k - \text{Exp}_k)
\end{equation}


\section{物理的描像：砂時計と金庫}
\begin{itemize}
    \item \textbf{期中（$t-1 \to t$）：} 売上と費用の砂時計（P/L）に砂が堆積する。
    \item \textbf{Closing（時刻 $t$）：} 砂時計の差分（利益）を測定し、会社の金庫（B/Sの $\text{Equity}$）に流し込む。その後、砂時計を反転・ゼロクリアして次期に備える。
\end{itemize}

P/Lは毎年破棄される揮発性メモリであり、B/Sは企業の全歴史を記録する不揮発性ストレージである。この境界線をリフレッシュする処理がClosingである。

\else
\chapter{Closing as a Discrete-Time Register Reset}
\label{chap:closing_model}

The mystified ritual of "closing accounts" is purely a discrete-time register clear and difference-accumulation cycle.

\section{Recurrence Formulation of Equity Evolution}
Let $t \in \{0, 1, 2, \dots\}$ denote discrete accounting periods. Net income over $(t-1, t]$ is defined as:
\begin{equation}
    \Delta \text{Equity}_t \equiv \text{Rev}_t - \text{Exp}_t
\end{equation}
\begin{equation}
    \mathbf{\text{Equity}_t = \text{Equity}_{t-1} + \Delta \text{Equity}_t}
\end{equation}
\begin{equation}
    \text{Equity}_n = \text{Equity}_0 + \sum_{k=1}^n (\text{Rev}_k - \text{Exp}_k)
\end{equation}


\section{Physical Representation: Hourglass and Vault}
\begin{itemize}
    \item \textbf{During Period ($t-1 \to t$):} Sand accumulates in the volatile revenue and expense hourglass (P/L).
    \item \textbf{Closing (Time $t$):} The difference (Net Income) is transferred to the non-volatile corporate vault (B/S Equity). The hourglass is zero-reset for the subsequent period.
\end{itemize}
\fi
